\begin{lemma}[Koll\'ar]
\label{lemma-hart-serre-loc-thm}
\begin{reference}
This is taken from a forthcoming paper by
J\'anos Koll\'ar entitled ``Variants of normality for Noetherian schemes''.
\end{reference}
Let $(R, \mathfrak m)$ be a local Noetherian ring.
Then exactly one of the following holds:
\begin{enumerate}
\item $(R, \mathfrak m)$ is Artinian,
\item $(R, \mathfrak m)$ is regular of dimension $1$,
\item $\text{depth}(R) \geq 2$, or
\item there exists a finite ring map $R \to R'$ which is not
an isomorphism whose kernel and cokernel are annihilated by a power
of $\mathfrak m$ such that $\mathfrak m$ is not an associated
prime of $R'$ and $R' \not = 0$.
\end{enumerate}
\end{lemma}

\begin{proof}
Observe that $(R, \mathfrak m)$ is not Artinian if and only if
$V(\mathfrak m) \subset \Spec(R)$ is nowhere dense. See
Proposition \ref{proposition-dimension-zero-ring}. We assume this from now on.

\medskip\noindent
Let $J \subset R$ be the largest ideal killed by a power of $\mathfrak m$.
If $J \not = 0$ then $R \to R/J$ shows that $(R, \mathfrak m)$
is as in (4).

\medskip\noindent
Otherwise $J = 0$. In particular $\mathfrak m$ is not an associated prime
of $R$ and we see that there is a nonzerodivisor $x \in \mathfrak m$ by
Lemma \ref{lemma-ideal-nonzerodivisor}. If $\mathfrak m$
is not an associated prime of $R/xR$ then $\text{depth}(R) \geq 2$
by the same lemma. Thus we are left with the case when there is a
$y \in R$, $y \not \in xR$ such that $y \mathfrak m \subset xR$.

\medskip\noindent
If $y \mathfrak m \subset x \mathfrak m$ then we can consider the
map $\varphi : \mathfrak m \to \mathfrak m$, $f \mapsto yf/x$
(well defined as $x$ is a nonzerodivisor). By the determinantal trick
of Lemma \ref{lemma-charpoly-module} there exists a monic
polynomial $P$ with coefficients in $R$ such that $P(\varphi) = 0$.
We conclude that $P(y/x) = 0$ in $R_x$.
Let $R' \subset R_x$ be the ring generated by
$R$ and $y/x$. Then $R \subset R'$ and $R'/R$ is a finite $R$-module
annihilated by a power of $\mathfrak m$. Thus $R$ is as in (4).

\medskip\noindent
Otherwise there is a $t \in \mathfrak m$ such that $y t = u x$
for some unit $u$ of $R$. After replacing $t$ by $u^{-1}t$
we get $yt = x$. In particular $y$ is a nonzerodivisor.
For any $t' \in \mathfrak m$ we have $y t' = x s$ for some $s \in R$.
Thus $y (t' - s t ) = x s - x s = 0$.
Since $y$ is not a zero-divisor this implies that $t' = ts$ and so
$\mathfrak m = (t)$. Thus $(R, \mathfrak m)$ is regular of dimension 1.

\medskip\noindent
The argument so far shows that every $R$ falls into one of the 4 cases.
To finish we have to show that no two of the properties can hold
simultaneously for each of the 6 combinations of properties. We'll
just show that (2) and (3) each exclude (4); the other cases are left
to the reader.

\medskip\noindent
Assume $R$ is regular of dimension $1$ and that $R \to R'$ is a finite
ring map whose kernel and cokernel are annihilated by a power of $\mathfrak m$
and $\mathfrak m$ is not an associated prime of $R'$. Then $R'$
is a finite $R$-module of depth at least $1$ and hence free by
Lemma \ref{lemma-regular-mcm-free}. Since $R \to R'$ is an isomorphism
in the generic point, we conclude that $R'$ must be free of rank $1$
as an $R$-module. Then $R' \to \text{End}_R(R') \cong R$ is an inverse
to the map $R \to R'$. Thus (4) cannot be true.

\medskip\noindent
Assume $R$ has depth $\geq 2$ and that $R \to R'$ is a finite
ring map whose kernel and cokernel are annihilated by a power
of $\mathfrak m$ and $\mathfrak m$ is not an associated prime of $R'$.
Then $R \to R'$ is necessarily injective (as the kernel would have
both depth $0$ and depth $\geq 1$) and the cokernel has depth
at least $1$ (by Lemma \ref{lemma-depth-in-ses}
and the fact that $R'$ has depth $\geq 1$ by assumption)
whence cannot be supported on
$\{\mathfrak m\}$ unless it is $0$ as well. Thus (4) cannot be true.
\end{proof}